\section{Validation}

The implementation is validated with automated tests based on JUnit 5 and the Pekko typed test kit.
The test suite covers:
\begin{itemize}
    \item configuration loading from \texttt{application.conf};
    \item domain-model validation for sensors, events, states, and zones;
    \item keypad forwarding of PIN and arming requests;
    \item sensor forwarding of activation events;
    \item siren activation, deactivation, and state queries;
    \item complete control-unit transitions across all alarm states;
    \item partial arming and ignored events from inactive zones;
    \item stale timeout messages received in states where they are no longer relevant;
    \item CLI command parsing and demo execution.
\end{itemize}

The most important tests are the control-unit tests, because the control unit is the only owner of the alarm state.
They verify that sensors are ignored in \textit{DISARMED} and \textit{EXIT\_DELAY}, that active-zone sensors start
\textit{ENTRY\_DELAY}, that the entry timeout activates the siren, and that the correct PIN always returns the system to
\textit{DISARMED}.

From a concurrency point of view, the relevant safety property is that no actor shares mutable state with another actor.
The control unit processes one message at a time from its mailbox, so concurrent inputs are serialized before they can
change the alarm behavior.
Timers are checked as ordinary messages, which makes timed transitions testable in the same way as keypad and sensor
events.

The resulting coverage validates both the functional requirements and the message-passing structure required by the
assignment.
